% =====================================================================
% model.tex
% Content: the signal-detection-theoretic model of the cued change-
% detection task. Task structure, attention allocation, attention-to-
% perception mapping with asymmetric benefit/cost, reward structure and
% the decorrelation channel (the correlated-noise orthant integral), and
% the four-policy decomposition that defines the three levers
% (criterion, sensitivity, decorrelation).
% =====================================================================

\section{Model}
\label{sec:model}

The model is a signal-detection-theoretic account of a cued
change-detection task with four nested decision policies. It admits
three mechanisms by which an observer can adapt to a value cue:
criterion adjustment, sensitivity re-allocation, and decorrelation of
cross-location noise (correlation $\corr$). When $\corr = 0$ the
decorrelation lever is inactive and the value response reduces to the
familiar criterion-versus-sensitivity description
(Section~\ref{sec:model-policies}). We first state the task and the
attention-to-perception machinery (Sections~\ref{sec:model-task}--%
\ref{sec:model-mapping}), then introduce the reward structure and the
single place the independence assumption enters it
(Section~\ref{sec:model-reward}), and finally the policy decomposition
that isolates each lever's contribution
(Section~\ref{sec:model-policies}).

% ---------------------------------------------------------------------
\subsection{Task structure}
\label{sec:model-task}

Consider a spatial change-detection task with $\Nloc \ge 2$ locations,
one of which is cued. Index the locations so that $i = c$ is the cued
location and $i = 1, \dots, \Nloc - 1$ enumerate the $\Nloc - 1$ uncued
locations. The cue conveys two pieces of information:
\begin{itemize}
\item \textbf{Value} ($\val \ge 1$): the reward for detecting a change
  at the cued location. Changes at uncued locations are worth a reward
  of $1$.
\item \textbf{Validity} ($\valid \in [1/\Nloc, 1]$): the probability
  that, given a change trial, the change occurs at the cued location.
  Each uncued location receives probability $(1 - \valid)/(\Nloc - 1)$.
\end{itemize}
On each trial a change occurs with probability $0.5$ (a ``change
trial'') or does not (a ``no-change trial''). Each location carries an
equal-variance Gaussian decision variable, and the observer makes a
detection decision at each location using signal detection theory. The
per-location hit rate and false-alarm rate are
\begin{equation}
  \HR_i \;=\; \Phinorm(\dprime_i/2 - c_i),
  \qquad
  \FAR_i \;=\; \Phinorm(-\dprime_i/2 - c_i),
  \label{eq:sdt-marginal}
\end{equation}
where $\dprime_i$ is perceptual sensitivity at $i$, $c_i$ is the
decision criterion at $i$, and $\Phinorm$ is the standard-normal
cumulative distribution function \cite{MullerFindlay1987,
posner1980_orienting}.

% ---------------------------------------------------------------------
\subsection{Attention allocation}
\label{sec:model-allocation}

The observer allocates a continuous share of attention
$\alphacued \in [0, 1]$ to the cued location. The remaining attention is
distributed equally among the uncued locations, so each receives
$\alphauncued = (1 - \alphacued)/(\Nloc - 1)$ and the shares satisfy
$\alphacued + (\Nloc - 1)\,\alphauncued = 1$. At uniform attention
$\alphacued = 1/\Nloc$ and all locations are treated identically.

% ---------------------------------------------------------------------
\subsection{Attention-to-perception mapping}
\label{sec:model-mapping}

Attention determines perceptual sensitivity at each location through a
baseline transfer function
\begin{equation}
  f(a) \;=\; f_0 + (1 - f_0)\,h(a),
  \label{eq:transfer}
\end{equation}
where $a$ is the attention allocated to a location, $f_0 \in (0, 1)$ is
the baseline sensitivity fraction at minimal attention, and $h$ is a
monotonically increasing function with $h(0) = 0$ and $h(1) = 1$. We
consider four functional forms,
\begin{equation}
  h(a) \;\in\; \{\, a,\ \sqrt{a},\ a^{0.3},\ a^{2} \,\},
  \label{eq:h-forms}
\end{equation}
spanning regimes from strongly diminishing returns ($a^{0.3}$) to
accelerating returns ($a^{2}$). At uniform attention every location
achieves a baseline sensitivity $\dprime_{\mathrm{base}} = \dprimemax
\cdot f(1/\Nloc)$.

\paragraph{Independent benefit and cost.}
Attentional re-allocation need not help and hurt equally. We
parameterise the asymmetry between attentional benefit and cost with a
ratio $\Rsens > 0$. When a location gains attention relative to uniform,
its $\dprime$ departure from baseline is scaled by a benefit weight
$\benefit(\Rsens)$; when a location loses attention, the departure is
scaled by a cost weight $\cost(\Rsens)$. Under an additive
conservation rule,
\begin{equation}
  \benefit(\Rsens) = \frac{2\Rsens}{\Rsens + 1}
  \quad\text{(benefit)},
  \qquad
  \cost(\Rsens) = \frac{2}{\Rsens + 1}
  \quad\text{(cost)},
  \label{eq:beta-gamma}
\end{equation}
which satisfy $\benefit + \cost = 2$ (conserving total magnitude) and
$\benefit/\cost = \Rsens$. The ratio $\Rsens$ reflects the relative
efficacy of attentional enhancement versus suppression: a
benefit-dominant regime ($\Rsens > 1$) corresponds to circuits in which
top-down feedback boosts attended representations more effectively than
lateral inhibition suppresses unattended ones, and a cost-dominant
regime ($\Rsens < 1$) to strong surround suppression at unattended
locations \cite{reynolds_heeger2009_normalization,
treue_martinez_trujillo1999_feature_attention,
mcadams_maunsell1999_v4_tuning}. The resulting per-location sensitivities
are
\begin{align}
  \dprime_c
    &= \dprime_{\mathrm{base}}
       + \benefit(\Rsens) \cdot
       \big[\, \dprimemax\, f(\alphacued) - \dprime_{\mathrm{base}} \,\big],
  \label{eq:dprime-cued}\\
  \dprime_u
    &= \dprime_{\mathrm{base}}
       + \cost(\Rsens) \cdot
       \big[\, \dprimemax\, f(\alphauncued) - \dprime_{\mathrm{base}} \,\big],
  \label{eq:dprime-uncued}
\end{align}
with the convention that if $\alphacued < 1/\Nloc$ the roles reverse
(the cued departure is scaled by $\cost$ and each uncued by $\benefit$)
and all $\dprime$ values are clamped at $\ge 0$. At $\Rsens = 1$,
$\benefit = \cost = 1$ and the model reduces to a symmetric special case
in which a single shared transfer function governs both benefit and
cost; this symmetric corner is the smooth centre of the asymmetry
family, and its exact consistency properties are recorded in the
Supplementary material (Section~\ref{sec:appendix}).

\textcolor{red}{[GAP G-001: needs the attention-to-$\dprime$ mapping
figure.]}

% ---------------------------------------------------------------------
\subsection{Reward structure and the decorrelation channel}
\label{sec:model-reward}

We evaluate two reward variants to separate task-design effects from
fundamental incentive structure. \emph{Variant~A} (value-coupled
correct rejection) couples the correct-rejection reward to value and
validity, $\CR = \valid\,\val + (1 - \valid)$, calibrated so that the
expected value of responding and withholding are balanced at uniform
attention. \emph{Variant~B} (fixed correct rejection) sets $\CR = 1$
independent of $\val$ and $\valid$, which creates a stronger marginal
benefit for detection at high-value locations because the hit--$\CR$
difference at the cued position scales with $\val$. The expected reward
on a single trial is
\begin{equation}
  \mathbb{E}[\Rew]
  \;=\;
  \tfrac{1}{2}\big[\,\valid\,\HR_c\,\val
                  \,+\,(1-\valid)\,\HR_u\,\big]
  \;+\;
  \tfrac{1}{2}\,\PnoFA\,\CR,
  \label{eq:expected-reward}
\end{equation}
where $\PnoFA$ is the probability of no false alarm at any location on a
no-change trial.

\paragraph{The locus of the independence assumption.}
Equation~(\ref{eq:expected-reward}) is linear in the marginal hit rates
$\HR_c, \HR_u$: on a change trial the change occurs at a single location,
so no product of probabilities across locations appears in the
change-trial bracket. The \emph{only} cross-location term in
(\ref{eq:expected-reward}) is $\PnoFA$. Under independent per-location
decision variables,
\begin{equation}
  \PnoFA^{\text{indep}}
  \;=\;
  \prod_{i} (1 - \FAR_i)
  \;=\;
  \Phinorm(b_c)\,\Phinorm(b_u)^{\Nloc - 1},
  \qquad
  b_i := c_i + \dprime_i/2,
  \label{eq:pnofa-indep}
\end{equation}
where $b_i$ is the $z$-score the no-change decision variable must stay
below. Equation~(\ref{eq:pnofa-indep}) is the sole place the
independence assumption enters the reward structure, so relaxing that
assumption amounts to replacing this product by the correlated joint
orthant probability.

\paragraph{The decorrelation channel.}
We promote the independence assumption to a tunable equicorrelation
parameter $\corr \in [0, 1)$. Let the no-change decision variables be
jointly Gaussian with standardised marginals and equicorrelation,
\begin{equation}
  \mathbf{X} \;\sim\; \mathcal{N}(\boldsymbol\mu, \Sigma),
  \qquad
  \mu_i = -\dprime_i/2,
  \qquad
  \Sigma_{ii} = 1,
  \qquad
  \Sigma_{ij} = \corr \ (i \neq j).
  \label{eq:equicorr-cov}
\end{equation}
The one-factor representation
$X_i = \mu_i + \sqrt{\corr}\,Z + \sqrt{1 - \corr}\,\varepsilon_i$, with
$Z, \varepsilon_1, \dots, \varepsilon_\Nloc$ i.i.d.\ $\mathcal{N}(0, 1)$,
reproduces (\ref{eq:equicorr-cov}) and makes the $X_i$ conditionally
independent given $Z$. Multiplying the conditional event probabilities
and integrating out $Z$ against the standard-normal density $\varphi$
yields the exact one-dimensional integral
\begin{equation}
  \boxed{\;
    \PnoFA(\corr)
    \;=\;
    \int_{-\infty}^{\infty}
    \Phinorm\!\left(\frac{b_c - \sqrt{\corr}\,z}{\sqrt{1 - \corr}}\right)
    \Phinorm\!\left(\frac{b_u - \sqrt{\corr}\,z}{\sqrt{1 - \corr}}\right)^{\!\Nloc - 1}
    \varphi(z)\,dz.
  \;}
  \label{eq:pnofa-rho}
\end{equation}
Equation~(\ref{eq:pnofa-rho}) is exact for any $\Nloc$ and any
$\corr \in [0, 1)$; equicorrelation collapses the $\Nloc$-dimensional
orthant integral to a one-dimensional one, with no multivariate-normal
CDF required. We evaluate it by $64$-node Gauss--Hermite quadrature,
with quadrature error $\le 10^{-15}$ against a $128$-node reference
across the operating band $\corr \in \{0, 0.05, 0.1, 0.2, 0.3, 0.4\}$.

\paragraph{Recovery contract.}
At $\corr = 0$ the integrand of (\ref{eq:pnofa-rho}) is constant in $z$
and the integral reduces to (\ref{eq:pnofa-indep}):
\begin{equation}
  \PnoFA(0) \;=\;
  \Phinorm(b_c)\,\Phinorm(b_u)^{\Nloc - 1}.
  \label{eq:rho-zero-recovery}
\end{equation}
Equation~(\ref{eq:rho-zero-recovery}) is the \emph{recovery contract}:
with $\corr$ clamped to zero, the model's reward, VDA, and
criterion-fraction values coincide to floating-point identity with the
independent-noise model, so every result reported below contains the
independent-noise case as its $\corr = 0$ limit. The empirical envelope
we explore is $\corr \in [0, 0.4]$, which brackets the $\corr \approx 0.2$
spike-count correlation reported in macaque V4 during peripheral-cue
change detection \cite{CohenMaunsell2009}; higher $\corr$ values are
admissible in the model but anchor no headline claim. Equicorrelation is
a deliberate one-parameter simplification; structured covariances
\cite{RuffCohen2016, Srinath2021} break the exact one-dimensional
reduction (\ref{eq:pnofa-rho}) and are noted as a scoped limitation in
the Discussion (Section~\ref{sec:discussion}). By the Slepian
inequality, $\PnoFA(\corr)$ is non-decreasing in $\corr$ at fixed
$(b_c, b_u)$ \cite{Slepian1962}; the full derivation of
(\ref{eq:pnofa-rho}) and this monotonicity is in the Supplementary
material (Section~\ref{sec:appendix}).

% ---------------------------------------------------------------------
\subsection{Policy decomposition}
\label{sec:model-policies}

To isolate the contribution of each mechanism, we define four nested
policies that optimise (\ref{eq:expected-reward}) over progressively
shrinking parameter sets:
\begin{enumerate}
\item \textbf{Joint optimum (P1)}: $\alphacued, c_c, c_u$ jointly
  optimised at each $\val$. This is the best any observer can achieve.
\item \textbf{Value-blind attention (P2)}: $\alphacued$ fixed at its
  $\val = 1$ optimum --- the optimal allocation when all locations have
  equal value --- with the criteria re-optimised at each $\val$. This
  observer uses attention for validity but not for value.
\item \textbf{Uniform attention (P3)}: $\alphacued = 1/\Nloc$, criteria
  optimised. This observer uses only the criterion to encode value.
\item \textbf{Floor (P4)}: $\alphacued = 1/\Nloc$, $c = 0$. An unbiased
  baseline that uses neither mechanism.
\end{enumerate}
The reward gain decomposes into three additive pieces,
\begin{align}
  \text{criterion gain} &= \Rpthree - \Rpfour,
  \label{eq:gain-criterion}\\
  \text{validity-attention gain} &= \Rptwo - \Rpthree,
  \label{eq:gain-validity}\\
  \VDA &= \Rpone - \Rptwo,
  \label{eq:vda-def}
\end{align}
where the value-directed-attention benefit $\VDA$ --- the reward gained
by making attention track cue \emph{value} on top of an already
optimal, validity-driven allocation and criterion --- is our central
metric. The criterion fraction,
\begin{equation}
  \CF \;=\;
  \frac{\Rpthree - \Rpfour}{\Rpone - \Rpfour},
  \label{eq:cf-def}
\end{equation}
measures the share of the total reward gain captured by the criterion
lever alone. We report the range and distribution of $\CF$ across the
full parameter space in Section~\ref{sec:results}.

\paragraph{Three levers.}
At fixed correlation the value response decomposes into two levers:
criterion shift (P3~$\to$~P4) and sensitivity re-allocation
(P1~$\to$~P2). With $\corr$ a free model parameter, the same
expected-reward expression (\ref{eq:expected-reward}) admits a third.

\begin{definition}[The three levers]
\label{def:three-levers}
At fixed $(\Nloc, \valid, \val, \dprimemax, f_0, h, \CR, \Rsens)$, the
model exposes three independent mechanisms by which an observer can
adapt to a non-trivial cue:
\begin{enumerate}
\item \textbf{Criterion shift} --- moving $(c_c, c_u)$ off the
  $\val = 1$ floor. Its standalone reward contribution is
  $\Rpthree - \Rpfour$, and its share of the value-related gain is
  $\CF$ (\ref{eq:cf-def}).
\item \textbf{Sensitivity re-allocation} --- moving $\alphacued$ off the
  $\val = 1$ optimum. Its standalone contribution at adapted criteria is
  $\Rpone - \Rptwo = \VDA$ (\ref{eq:vda-def}).
\item \textbf{Decorrelation} --- moving $\corr$ off the
  $\corr = 0$ corner. By Slepian monotonicity $\PnoFA(\corr)$ is
  non-decreasing in $\corr$ at fixed $(b_c, b_u)$ \cite{Slepian1962},
  which lifts every per-policy supremum reward
  $\Rew_{\mathrm{P}_k}(\corr)$ above its $\corr = 0$ value.
\end{enumerate}
\end{definition}

The decorrelation lever is grounded in physiology: attention is known to
reduce interneuronal correlations in visual cortex
\cite{CohenMaunsell2009, RuffCohen2016, Srinath2021}. Treating $\corr$
as a model parameter rather than fixing it at zero lets us report what
the expected-reward expression does as cross-location noise is varied.
Whether independent-noise analyses bound the VDA benefit, and what
quantity they bound, is an empirical question about the sign of
$\partial\VDA/\partial\corr$ that we take up with the corresponding
sweep in the Results (Section~\ref{sec:results}) and Discussion
(Section~\ref{sec:discussion}).
